\documentclass[11pt]{article}
\usepackage[letterpaper,margin=1in]{geometry}
\usepackage{amsmath,amssymb,amsthm,mathtools,booktabs,microtype}
\usepackage[T1]{fontenc}
\usepackage{lmodern}
\usepackage[hidelinks]{hyperref}
\allowdisplaybreaks[2]
\newtheorem{theorem}{Theorem}[section]
\newtheorem{lemma}[theorem]{Lemma}
\newtheorem{corollary}[theorem]{Corollary}
\newtheorem{proposition}[theorem]{Proposition}
\theoremstyle{definition}\newtheorem{remark}[theorem]{Remark}
\newcommand{\C}{\mathbb C}
\newcommand{\E}{\mathbb E}
\newcommand{\Z}{\mathbb Z}
\newcommand{\supp}{\operatorname{supp}}
\newcommand{\ind}{\mathbf 1}
\newcommand{\abs}[1]{\lvert #1\rvert}
\newcommand{\norm}[1]{\lVert #1\rVert}
\newcommand{\Gq}{G_q}
\newcommand{\GRH}{\textup{DGRH}}
\title{Sparse Dirichlet mixtures as error-correcting codes:\\prime sampling and rigidity of near-minimum support}
\author{Research note for the ProveIt project}
\date{27 November 2026}
\begin{document}
\maketitle
\begin{abstract}
Assume the Riemann hypothesis and the generalized Riemann hypothesis for Dirichlet $L$-functions, as exported by the pinned Lean development \texttt{grh-lean}. For $2\le s\le Q$ we construct a single set of $O(s^6\log Q)$ distinct primes of size $O(s^6\log^2 Q)$ on which every $s$-term sum of characters modulo any $q\le Q$ has near-optimal support. The main new ingredient is a finite-group rigidity theorem valid for nonuniform sampling measures: if a normalized $s$-sparse character sum exceeds $1/(100s)$ on at most $(1+1/(100s^2))/s$ of the sample, then its frequencies are \emph{exactly} a coset of a subgroup of order $s$, and the function is uniformly close to a modulated coset indicator. At zero tolerance this identifies all near-minimum words and gives an exact formula for the minimum Hamming distance of the nonlinear $k$-sparse evaluation code whenever $2k\mid\varphi(q)$. We obtain the optimal limiting adversarial error fraction $1/(4k)$, a stability estimate for noisy data, and a mixed-conductor version requiring only small bias for characters of moduli at most $Q^4$. A cheaper $O(s^2\log Q)$ sampling regime gives near-optimal recovery without the rigidity conclusion. The analytic input is classical; the quantitative rigidity and its prime-code applications are the proposed contributions. No additional unproved hypothesis is used.
\end{abstract}
\tableofcontents
\section{Input, scope and main statements}
\paragraph{Upstream input.} We take the five theorem statements of the readable edition of the \texttt{grh-lean} repository, commit \texttt{72fead10} (full hash in the bibliography), as the input theorem. This is the edition vendored in ProveIt \cite{vendor,upstream}. In particular, the declarations \texttt{LFunction\_eq\_zero\_re} and \texttt{riemannHypothesis} provide exactly the critical-line statements needed below. We write \GRH{} for these inputs together. This is Dirichlet GRH, not GRH for arbitrary Artin, Hecke or automorphic $L$-functions. Independent machine checks at the pinned commit have been reported \cite{sutherland}; this note has not rerun them or formalized its own arguments. All results below using \GRH{} are implications from that theorem, and become unconditional if the input is accepted.

For an integer $q\ge2$ let $G_q=(\Z/q\Z)^\times$ and let $\widehat G_q$ be its character group. A character sum is a function $F:G_q\to\C$ with Fourier expansion $F=\sum_\chi c_\chi\chi$. Its sparsity is $\abs{\supp c}$, and its coefficient norm is $\norm{c}_2=(\sum\abs{c_\chi}^2)^{1/2}$. For a finite set $S$ of primes, put $S_q=\{p\in S:p\nmid q\}$ and $n_q=\abs{S_q}$. Sample points are counted with multiplicity as residues, but the primes themselves are distinct.

Here and below all logarithms are natural. Fix an integer $s\ge2$ and put
\begin{equation}\label{eq:constants}
 \eta=\frac1{1000s^3},\qquad \epsilon=\frac1{100s^2},\qquad \tau_0=\frac1{100s}.
\end{equation}
For $0\le \tau\le\tau_0$ define the lower large-value fraction
\begin{equation}\label{eq:lower}
 d_s(\tau)=\frac{1-(s-1)\eta-\tau^2}{s}.
\end{equation}

\begin{theorem}[Universal prime rigidity, using \GRH]\label{thm:main}
There are absolute constants $B,K>0$ such that for every integer $Q\ge2$ and $2\le s\le Q$ there is a set $S$ of distinct primes contained in $[X,2X]$, with
\[
 X=K\eta^{-2}\log^2(2Q),\qquad
 B\eta^{-2}\log(2Q)/2\le\abs S\le 2B\eta^{-2}\log(2Q),
\]
with the following simultaneous property for all $2\le q\le Q$. Every nonzero $F$ with at most $s$ Fourier terms satisfies
\begin{equation}\label{eq:largeset}
 \frac{\#\{p\in S_q:\abs{F(p)}>\tau\norm c_2\}}{n_q}\ge d_s(\tau)
 \qquad(0\le\tau\le\tau_0).
\end{equation}
If the left side is at most $(1+\epsilon)/s$, then its Fourier support is exactly $\chi_0H$ for a subgroup $H\le\widehat G_q$ of order $s$. Write $H^\perp=\{u\in G_q:h(u)=1\text{ for every }h\in H\}$. There is a coset $C=gH^\perp$ and a nonzero $\lambda\in\C$ such that
\begin{equation}\label{eq:stability}
 \norm{F-\lambda\chi_0\ind_C}_{\infty,G_q}\le\tau\norm c_2,
 \qquad
 \norm{c-c_*}_2\le \tau\sqrt{1-1/s}\,\norm c_2,
\end{equation}
where $c_*$ is the coefficient vector of $\lambda\chi_0\ind_C$. Moreover $F=\lambda\chi_0$ on $C$, and $\abs{F}>\tau\norm c_2$ holds precisely on $C$, on the entire group.
\end{theorem}

The conclusion is about the exact algebraic support, not just a vector close to some sparse vector. For $\tau=0$ the only functions with sufficiently small empirical support are the modulated coset indicators, even though the sample need not contain every element of $G_q$. Conversely all such indicators lie within the near-minimum window because their sampled masses differ from $1/s$ by at most $(s-1)\eta/s$.

\begin{theorem}[Exact distance and recovery, using \GRH]\label{thm:code}
Let $s=2k\le Q$, let $S$ be as in Theorem~\ref{thm:main}, and for each $q\le Q$ form the nonlinear code
\[
 \mathcal C_{q,k}=\{(F(p))_{p\in S_q}:\abs{\supp c}\le k\}\subseteq\C^{n_q}.
\]
Its minimum Hamming distance $D_{q,k}$ is positive. If $s\mid\varphi(q)$, then exactly
\begin{equation}\label{eq:exactdist}
 D_{q,k}=\min_{\substack{H\le\widehat G_q,\ |H|=s\\gH^\perp\in G_q/H^\perp}}
     \#\{p\in S_q:p\bmod q\in gH^\perp\}.
\end{equation}
Every difference attaining this distance is a modulated coset indicator with a minimizing coset; conversely all such indicators are differences of two $k$-sparse sums. If $s\nmid\varphi(q)$, then $D_{q,k}>(1+\epsilon)n_q/s$.

For exact data, the optimal worst-case number of corruptions correctable by unique decoding is $\lfloor(D_{q,k}-1)/2\rfloor$. For noisy data, suppose $y$ agrees with each of two $k$-sparse evaluation vectors to accuracy $v\ge0$ outside at most $e$ positions for each vector. If $2e<d_s(\tau_0)n_q$, then their coefficient vectors $c,c'$ satisfy
\begin{equation}\label{eq:noise}
 \norm{c-c'}_2\le 200s\,v.
\end{equation}
If $2e\le(1+\epsilon)n_q/s$ and this norm exceeds $200s\,v$, then the normalized difference has the exact subgroup-frequency structure and the approximation in Theorem~\ref{thm:main}. In particular the two coefficient supports are disjoint and each has exactly $k$ elements.
\end{theorem}
The fraction of errors correctable is asymptotically $1/(4k)$ as the bias is made arbitrarily small (or, in the regime above, as $k$ grows), whenever the subgroup obstruction exists. More precisely the normalized minimum distance in \eqref{eq:exactdist} lies in $[(1-(s-1)\eta)/s,(1+(s-1)\eta)/s]$, and the exact integer radius is given above. In fact the upper bound can be sharpened to $D_{q,k}\le\lfloor n_q/s\rfloor$: the $s$ cosets for any one subgroup partition the sample, so at least one has count at most the average. For a fixed $k$, any smaller fixed positive bias can be used for all sufficiently large $Q$ by Lemma~\ref{lem:sample}. Hence no universal guarantee above the error fraction $1/(4k)$ is possible in these moduli, for \emph{any} prime sample, not just uniform ones. Here and in what follows uniqueness refers to the information-theoretic problem, not polynomial-time decoding.

\section{The analytic supply of small-bias primes}
\label{sec:analytic}
We give the analytic input in the form useful for this application. It is the same method as the merged article's smoothed prime estimate, prime Gram estimate and universal prime-set lemma \cite[labels \texttt{dir:lem:smooth}, \texttt{rec:lem:gram}, \texttt{dir:lem:oneset}]{article}. We need the dependence on a bias parameter that can decrease with $Q$.

\begin{lemma}[Uniform smoothed prime estimate, using \GRH]\label{lem:prime}
Fix a nonnegative $C^\infty$ function $W$ supported in $(1,2)$ with $\int W(u)\,du=1$. For every $x\ge2$ and every Dirichlet character $\chi$ modulo $h\ge1$, extended by zero off the units,
\begin{equation}\label{eq:prime}
 \sum_p(\log p)W(p/x)\chi(p)
 =\ind_{\chi=\chi_{0,h}}x+O_W(\sqrt x\log(2h)).
\end{equation}
The constant is independent of $x$ and $h$.
\end{lemma}
\begin{proof}
First let $\psi$ be primitive of conductor $f$ (including the principal character of conductor $1$, whose function is $\zeta$). For $\widehat W(z)=\int_0^\infty W(u)u^{z-1}\,du$, Mellin inversion and the Euler product on $\Re z=2$ give
\[
 \sum_{n\ge1}\Lambda(n)W(n/x)\psi(n)
 =\frac1{2\pi i}\int_{2-i\infty}^{2+i\infty}
       -\frac{L'}{L}(z,\psi)\,\widehat W(z)x^z\,dz.
\]
On $-1/2\le\Re z\le2$, repeated integration by parts gives $\widehat W(z)=O_{W,N}((1+|\Im z|)^{-N})$ for any fixed $N$. Move the line to $-1/2$. The standard local zero-counting estimate for primitive Dirichlet functions gives $O(\log(2f(j+2)))$ zeros, with multiplicity, at $j\le |\Im\rho|<j+1$. The usual logarithmic-derivative partial fraction estimate allows horizontal lines tending to infinity at distance $\gg1/\log(2f(T+2))$ from the zeros, on which the logarithmic derivative is $O(\log^2(2f(T+2)))$ in this fixed strip (see \cite[Chs. 12--14]{MV} or \cite{Davenport}). Thus the horizontal integrals vanish. On $\Re z=-1/2$ the functional equation, Stirling's estimate and the absolutely convergent logarithmic derivative on $\Re(1-z)=3/2$ give $O(\log(2f(|\Im z|+2)))$. The integral on the shifted line is $O_W(x^{-1/2}\log(2f))$.

The pole of $\zeta$ contributes $x\widehat W(1)=x$; a possible trivial zero at zero contributes $O_W(1)$. The nontrivial zeros contribute $-\sum_\rho x^\rho\widehat W(\rho)$. By \GRH{} and the unit-band count this is absolutely convergent with absolute value $O_W(\sqrt x\log(2f))$. The half-plane theorem and functional equation exclude nontrivial zeros on the two strip boundaries $\Re z=1,0$; the possible trivial zero at $0$ was accounted for. There are no other poles or zeros in the shifted strip. This proves the desired estimate with the von Mangoldt weight for the primitive character.

Pass to primes: the total higher-prime-power contribution is bounded using Chebyshev's bound $\vartheta(t)\ll t$ by
\[
 \norm W_\infty\sum_{2\le m\le\log(2x)/\log2}\vartheta((2x)^{1/m})
 \ll_W \sqrt x+x^{1/3}\log(2x)\ll_W\sqrt x.
\]
Finally the primitive character inducing $\chi$ differs on prime values only at divisors of $h$; their total logarithmic weight is at most $\log h$. Its inducing character is principal exactly when $\chi$ is principal. Since $f\le h$, this proves \eqref{eq:prime}.
\end{proof}

\begin{lemma}[A universal distinct prime sample, using \GRH]\label{lem:sample}
For every integer $Q\ge2$, write $L=\log(2Q)$. Let
\[
 (1000Q^3)^{-1}\le b\le1/4.
\]
There are absolute constants $B,K>0$, uniform in $b,Q$, for which there is a set $S$ of distinct primes in $[X,2X]$, $X=Kb^{-2}L^2$, such that, with $m=Bb^{-2}L$ and $n=|S|$,
\begin{align}
 m/2\le n\le 2m,\qquad
 \abs{\sum_{p\in S}\xi(p)}\le bn/8
 \quad\text{for all nonprincipal }\xi\bmod h,\ 2\le h\le Q^4.\label{eq:sample}
\end{align}
For every product $M$ of at most $s\le Q$ integers at most $Q$, provided $sb\le1/4$, the set $S_M=\{p\in S:p\nmid M\}$ obeys
\begin{equation}\label{eq:remove}
 |S\setminus S_M|\le bn/8.
\end{equation}
If a nonprincipal character $\xi$ modulo $h\mid M$ has $h\le Q^4$, its pullback $\widetilde\xi$ to the units modulo $M$ has
\begin{equation}\label{eq:conditionalbias}
 \left|\frac1{|S_M|}\sum_{p\in S_M}\widetilde\xi(p)\right|\le b.
\end{equation}
The same conclusions hold if $M$ is the least common multiple of these integers.
\end{lemma}
\begin{proof}
Set $w_p=(\log p)W(p/X)$ and $W_0=\sum_p w_p$. Lemma~\ref{lem:prime} with $h=1$ gives $W_0\ge X/2$ once $K$ is large. The same lemma gives, for all the nonprincipal characters under consideration,
\[
 \abs{\sum_p w_p\xi(p)}/W_0\ll_W L/\sqrt X\le b/32
\]
by increasing $K$. The implied constant is uniform because $\log(2h)\le4L$.

Select each prime independently with probability $m w_p/W_0$. These probabilities are at most one: since $b^{-2}\le 10^6Q^6$,
\[
 \frac{m w_p}{W_0}\le\frac{2B\norm W_\infty\log(2X)}{KL},\qquad
 \frac{\log(2X)}{L}\le C_1+C_2\log K
\]
for absolute $C_1,C_2$ and $K\ge1$. After choosing $B$, make $K$ sufficiently large to meet both conditions. The expected sample size is $m$. For any admissible $\xi$, the random sum $Z_\xi=\sum_{p\in S}\xi(p)$ has $|\E Z_\xi|\le bm/32$. The real and imaginary parts of its centered summands have absolute value at most $1$ and total variance at most $m$. Bernstein's inequality yields
\[
 \Pr(|Z_\xi-\E Z_\xi|>bm/32)\le4\exp(-c b^2m),
 \qquad \Pr(|n-m|>m/2)\le2\exp(-c' m)
\]
with absolute positive constants. (For the complex bound, apply the real inequality at threshold $bm/(32\sqrt2)$ to each part.) The number of characters with moduli at most $Q^4$ is at most $\sum_{h\le Q^4}h\le Q^8$. Choose $B$ so that the union of the failure probabilities is strictly less than one for every $Q\ge2$. A successful sample has $|Z_\xi|\le bm/16\le bn/8$ and $m/2\le n\le 3m/2$, proving \eqref{eq:sample}.

Because all sampled primes are at least $X$ and $M\le Q^s$, at most $s\log Q/\log X$ of them divide $M$. Dividing by $n\ge Bb^{-2}L/2$ gives a discarded fraction at most $2sb^2/(B\log X)\le b/8$ for large absolute $B$ and $K$ (use $sb\le1/4$). On $S_M$ the pulled-back character equals $\xi$ on integers. Removing the other primes changes the sum by at most their number. Consequently
\[
 \abs{\sum_{p\in S_M}\widetilde\xi(p)}/|S_M|
 \le\frac{b/8+b/8}{1-b/8}<b.
\]
All these estimates also apply to the least common multiple. The same sample works for every $M$; there is no union bound over such products.
\end{proof}

\section{An unconditional small-bias rigidity theorem}
\label{sec:rigid}
We now isolate the finite-dimensional contribution. Let $G$ be a finite abelian group written multiplicatively, and let $\mu$ be any probability distribution on $G$. Write $\widehat\mu(\gamma)=\E_\mu\gamma(g)$; our sign convention is immaterial because only absolute values are used. Let $A$ be a nonempty set of characters and $D=AA^{-1}$ be its ratio set. Assume the following \emph{local bias condition}:
\begin{equation}\label{eq:localbias}
 |\widehat\mu(\gamma)|\le\eta\quad\text{for every }1\ne\gamma\in DD^{-1}.
\end{equation}
It is weaker than small bias on all nonprincipal characters. It involves only products of at most four members of $A$ and their inverses.

\begin{lemma}[Gram bound and restricted Bessel inequality]\label{lem:gram}
If $\gamma_1,\dots,\gamma_r$ are distinct characters and $|\widehat\mu(\gamma_i\gamma_j^{-1})|\le b$ for $i\ne j$, then for $a\in\C^r$
\begin{equation}\label{eq:gram}
 (1-(r-1)b)\norm a_2^2\le\E_\mu|\sum_i a_i\gamma_i|^2
 \le(1+(r-1)b)\norm a_2^2.
\end{equation}
For a set $U\subseteq G$ of measure $\alpha>0$ and conditional measure $\mu_U$,
\begin{equation}\label{eq:bessel}
 \alpha\sum_{i=1}^r|\E_{\mu_U}\gamma_i|^2\le1+(r-1)b.
\end{equation}
\end{lemma}
\begin{proof}
Expand the square; the total absolute contribution of off-diagonal entries is at most $b((\sum |a_i|)^2-\sum |a_i|^2)\le b(r-1)\norm a_2^2$. For the second inequality the linear map $T:a\mapsto\sum a_i\gamma_i$ has squared operator norm at most $1+(r-1)b$ into $L^2(\mu)$. Apply the adjoint to $\ind_U$ of squared norm $\alpha$: $\sum |\E_\mu \ind_U\overline{\gamma_i}|^2\le(1+(r-1)b)\alpha$. Divide by $\alpha$.
\end{proof}

\begin{theorem}[Local-bias support rigidity, unconditional]\label{thm:finite}
Let $s\ge2$, $\eta,\epsilon,\tau_0$ be as in \eqref{eq:constants}. Suppose $F=\sum_{\chi\in A}c_\chi\chi$, all $c_\chi\ne0$, $t=|A|\le s$, and \eqref{eq:localbias} holds. For $0\le\tau\le\tau_0$, set
\[
 U=\{g\in G:|F(g)|>\tau\norm c_2\},\qquad\alpha=\mu(U).
\]
Then $\alpha\ge d_s(\tau)$. If $\alpha\le(1+\epsilon)/s$, necessarily $t=s$ and $A=\chi_0H$ with $H$ a subgroup of $\widehat G$ of order $s$. The conclusions of \eqref{eq:stability} hold, with $G$ in place of $G_q$. Conversely for $\tau=0$, every modulated nonzero indicator of a coset of $H^\perp$ has Fourier support $\chi_0H$ and satisfies the upper threshold if \eqref{eq:localbias} holds.
\end{theorem}
\begin{proof}
By homogeneity normalize $\norm c_2=1$. Since $1\in D$, all nontrivial ratios of members of $A$ belong to $DD^{-1}$. Lemma~\ref{lem:gram} gives $\E_\mu|F|^2\ge1-(t-1)\eta$. Pointwise $|F|^2\le t$ by Cauchy--Schwarz and on $G\setminus U$ it is at most $\tau^2$. Therefore
\begin{equation}\label{eq:basic}
 \alpha\ge\frac{1-(t-1)\eta-\tau^2}{t}\ge d_s(\tau)>0.
\end{equation}
Assume the upper threshold. Put $a=(s-1)\eta+\tau^2$. Our constants imply $a\le1/(500s^2)$. If $t\le s-1$, then \eqref{eq:basic} would give $(1-a)/(s-1)\le(1+\epsilon)/s$, impossible since
\[
 s(1-a)-(s-1)(1+\epsilon)\ge1-\frac1{500s}-\frac{s-1}{100s^2}>0.
\]
Hence $t=s$.

\paragraph{Step 1: near equality forces alignment on the observed support.}
Let $u_i(g)=c_i\chi_i(g)$ ($1\le i\le s$), let $\overline u(g)=F(g)/s$ denote their arithmetic mean, and use $\E_U$ for conditional expectation on $U$. Define the alignment defect
\[
 R=s-\E_U|F|^2\ge0.
\]
Since $\E_\mu|F|^2\le \alpha\E_U|F|^2+\tau^2$, the energy lower bound and the upper threshold imply
\begin{equation}\label{eq:defect}
 R\le s\frac{\epsilon+a}{1+\epsilon}
 \le s(\epsilon+(s-1)\eta+\tau^2)\le\frac1{80s}.
\end{equation}
The exact identity $s\sum_i|u_i-\overline u|^2=s\sum_i|u_i|^2-|F|^2=s-|F|^2$ holds at each $g$. For every pair $i,j$ it follows that
\begin{equation}\label{eq:pair}
 \E_U|u_i-u_j|^2\le2\E_U(|u_i-\overline u|^2+|u_j-\overline u|^2)\le 2R/s.
\end{equation}
At least one $|c_j|\ge1/\sqrt s$, and $||c_i|-|c_j||^2\le\E_U|u_i-u_j|^2$. Since $R\le1/8$, we have $|c_i|\ge1/(2\sqrt s)$ for every $i$. The expansion of \eqref{eq:pair} gives
\[
 2|c_i||c_j|\bigl(1-|\E_U(\chi_i\chi_j^{-1})|\bigr)\le2R/s.
\]
Consequently \emph{every} ratio in $D$ has conditional bias at least $1-4R\ge1-1/(20s)$.

\paragraph{Step 2: too many ratios contradict approximate orthogonality.}
If $|D|\ge s+1$, choose $s+1$ distinct ratios. By \eqref{eq:localbias} they obey the hypotheses of Lemma~\ref{lem:gram}. Inequalities \eqref{eq:bessel}, \eqref{eq:basic} and \eqref{eq:defect} now give
\begin{equation}\label{eq:contra}
 1+s\eta\ge(s+1)\alpha(1-1/(20s))^2
 \ge (1+1/s)(1-1/(500s^2))(1-1/(20s))^2.
\end{equation}
But the last expression is at least $(1+1/s)(1-1/(9s))$, by $(1-z)^2\ge1-2z$ and $(1-x)(1-y)\ge1-x-y$ for nonnegative $x,y\le1$. This is at least $1+5/(6s)$ for $s\ge2$, greater than $1+s\eta$. Contradiction. Thus $|D|\le s$. For any $\chi_0\in A$, $A\chi_0^{-1}\subseteq D$ has $s$ elements, so $D=A\chi_0^{-1}$. For any $h,h'\in D=A\chi_0^{-1}$ their ratio lies in $AA^{-1}=D$. Hence $D$ is a subgroup $H$ of order $s$, and $A=\chi_0H$.

\paragraph{Step 3: the exceptional set is exactly a coset.}
Let $K=H^\perp$; finite abelian duality gives $[G:K]=s$. For every coset $C=gK$ the exact character identity
\begin{equation}\label{eq:indicator}
 \ind_C(x)=\frac1s\sum_{h\in H}\overline{h(g)}h(x)
\end{equation}
holds (orthogonality of characters on $H$). Every $h\in H=D$ belongs to $DD^{-1}$, so
\begin{equation}\label{eq:mass}
 \frac{1-(s-1)\eta}{s}\le\mu(C)\le\frac{1+(s-1)\eta}{s}.
\end{equation}
Now $F/\chi_0$ is constant on each coset of $K$, hence $U$ is a union of cosets. It is nonempty by \eqref{eq:basic}; two cosets would already have measure at least $2(1-(s-1)\eta)/s>(1+\epsilon)/s$. Thus $U=C$ for exactly one coset. Let $\lambda=(F/\chi_0)|_C$; then $|\lambda|>\tau$, the difference $F-\lambda\chi_0\ind_C$ vanishes on $C$ and is bounded by $\tau$ elsewhere, on \emph{all} of $G$. Parseval with the uniform (not sampled) probability measure gives the coefficient norm bound $\tau\sqrt{1-1/s}$. Undo the normalization. Conversely \eqref{eq:indicator} has precisely $s$ nonzero Fourier coefficients after modulation, and \eqref{eq:mass} verifies the upper threshold since $(s-1)\eta\le\epsilon$.
\end{proof}

\begin{remark}[Relation to uncertainty principles]
For the uniform measure and $\tau=0$, \eqref{eq:basic} reduces to the classical Donoho--Stark support uncertainty principle, whose equality classification is known \cite{MOP,Meshulam}. The content here is the explicit nonzero gap for nonuniform small-bias observations and the approximate-zero conclusion: exact algebraic structure follows from a noisy empirical support bound. We do not claim optimality of the numerical rigidity gap, the bias exponent $s^{-3}$, or the sample complexity $s^6$. The local condition \eqref{eq:localbias} rather than a bound over all characters is important for the mixed-conductor application.
\end{remark}

\begin{proposition}[Generalized support weights, unconditional]\label{prop:rank}
Let $A$ have size $t$, and assume the bias bound $b$ on $AA^{-1}\setminus\{1\}$. Let $V\subseteq\C^A$ be an $r$-dimensional subspace of coefficient vectors ($1\le r\le t$), and let $F_1,\ldots,F_r$ arise from an orthonormal basis of $V$. For $u\ge0$,
\[
 \mu\{g:\sum_{i=1}^r |F_i(g)|^2>u^2\}
 \ge\frac{r(1-(t-1)b)-u^2}{t}.
\]
In particular the simultaneous support of $V$ has measure at least $r(1-(t-1)b)/t$.
\end{proposition}
\begin{proof}
By Lemma~\ref{lem:gram} the expected sum of squares is at least $r(1-(t-1)b)$. At each point the squared norm of the evaluation vector in $\C^A$ is $t$. Orthogonal projection onto the (conjugate) coefficient subspace gives $\sum |F_i(g)|^2\le t$ by Bessel. On the complement of the indicated event it is at most $u^2$. The claimed lower bound follows.
\end{proof}

\section{Applications to prime codes}
\label{sec:codes}
\begin{proof}[Proof of Theorem~\ref{thm:main}]
Apply Lemma~\ref{lem:sample} with $b=\eta$ of \eqref{eq:constants}. Its parameter range follows from $2\le s\le Q$. For each $q\le Q$, apply \eqref{eq:conditionalbias} with $M=h=q$ for nonprincipal characters. The empirical distribution of the residues of $S_q$ has bias at most $\eta$ on every nonprincipal character. In particular it satisfies \eqref{eq:localbias} for every Fourier support $A$. Now Theorem~\ref{thm:finite} gives the conclusions, simultaneously for all $q$ and $F$.
\end{proof}

\begin{proof}[Proof of Theorem~\ref{thm:code}]
The difference of two $k$-sparse sums has at most $s=2k$ terms. Conversely every sum with at most $s$ terms can be written as a difference of two sums with at most $k$ terms: partition its terms into two sets and negate the second set's coefficients. Theorem~\ref{thm:main} with $\tau=0$ implies injectivity and gives a positive lower bound for the Hamming weight of every nonzero difference. The minimum is attained since weights are integers in a finite range.

If $s\mid|G_q|$, a finite abelian group of order $|G_q|$ has a subgroup of order $s$ (e.g. use its primary decomposition); this applies to $\widehat G_q$. Every modulated coset indicator then has weight in the interval in \eqref{eq:mass}, multiplied by $n_q$. This upper bound is less than $(1+\epsilon)n_q/s$. Every minimum-weight difference must therefore have the exact shape in Theorem~\ref{thm:main} with $\tau=0$, proving \eqref{eq:exactdist} and the classification. If $s\nmid|G_q|$ no subgroup of order $s$ exists, so the rigidity theorem gives the strict gap instead.

For the decoding radius: two radius-$e$ Hamming balls are disjoint if $2e<D_{q,k}$ by the triangle inequality. Conversely, if $2e\ge D_{q,k}$, choose a pair at minimum distance and partition their differing positions into two sets of sizes at most $e$. A received word that copies each of the pair on the appropriate set lies in both balls. This proves the exact integer radius, irrespective of alphabet size.

For noisy data write $F$ for the difference, with coefficient vector $c-c'$. Outside the union of the two exceptional sets, $|F(p)|\le2v$. If $\norm{c-c'}_2>2v/\tau_0=200s v$, the event $|F|>\tau_0\norm{c-c'}_2$ is contained in a set of size at most $2e$. The lower bound \eqref{eq:largeset} contradicts $2e<d_s(\tau_0)n_q$. This proves \eqref{eq:noise}. Under the weaker upper bound $2e\le (1+\epsilon)n_q/s$, the same event satisfies the rigidity theorem. Its support contains exactly $s=2k$ characters, which forces each candidate to have exactly $k$ distinct terms with disjoint supports. This also proves the final assertion.
\end{proof}

\paragraph{Erasure and noise variants.} If $u$ erased coordinates are known, replace $2e$ everywhere by $2e+u$: the common good set loses at most this many coordinates. An exact word is recoverable precisely up to the worst-case condition $2e+u<D_{q,k}$ for given $e,u$ (when $2e+u\ge D$, partition the differing coordinates into up to $u$ erased and two groups of size at most $e$). The same proof gives \eqref{eq:noise} for the corresponding noisy model. The basic Gram bound also shows that every set of $t\le s$ columns has normalized Gram spectrum in $[1-(t-1)\eta,1+(t-1)\eta]$. Proposition~\ref{prop:rank} gives the displayed generalized support weights for this prime sample.

\section{Unknown and mixed conductors}
\label{sec:mixed}
A key distinction from using a single fixed group is that the common conductor of $s$ primitive characters of conductor at most $Q$ can be as large as $Q^s$. The sample size need not have this factor in its logarithm. We use the local rather than global bias condition.

Let $\mathcal P_Q$ be the finite dictionary of all primitive Dirichlet characters of conductor at most $Q$, including the trivial primitive character of conductor $1$. They are evaluated at integers with the usual zero convention at ramified primes. Define a mixed signal $F(n)=\sum_{\chi\in A}c_\chi\chi(n)$ for distinct $\chi\in\mathcal P_Q$. Its coefficient vector is indexed by $\mathcal P_Q$; unlike the previous section, the modulus is not part of the supplied data.

\begin{theorem}[Universal mixed-conductor recovery, using \GRH]\label{thm:mixed}
For the same $S$ as in Theorem~\ref{thm:main}, every nonzero mixed signal with $1\le|A|\le s$ satisfies
\begin{equation}\label{eq:mixedlower}
 \frac{\#\{p\in S:|F(p)|>\tau\norm c_2\}}{|S|}
 \ge (1-\eta/8)d_s(\tau),\qquad 0\le\tau\le\tau_0.
\end{equation}
Let $M$ be the least common multiple of the conductors of the terms. If the left side of \eqref{eq:mixedlower} is at most $(1-\eta/8)(1+\epsilon)/s$, the characters inflated to $G_M$ have support exactly a coset of a subgroup of size $s$; the function on $G_M$ has the stability and exact large-value coset conclusion of Theorem~\ref{thm:finite}. No classification of the additional ramified coordinates is asserted.

In particular, for $s=2k$ all $k$-sparse mixtures from the dictionary are uniquely identifiable from the common prime sample, and in the noisy two-candidate model of Theorem~\ref{thm:code} the bound \eqref{eq:noise} holds if
\[
 2e+u<(1-\eta/8)d_s(\tau_0)|S|,
\]
where $u$ denotes known erasures. This conclusion needs no advance knowledge of the conductors.
\end{theorem}
\begin{proof}
Inflate each primitive character to $G_M$. They remain distinct by uniqueness of the inducing primitive character. If $M=1$, the dictionary support is a singleton and $F$ is constant; the conclusion is immediate, with no near-minimum case. Otherwise write $A_M$ for the inflated set and $D=A_MA_M^{-1}$. Each nontrivial member $\gamma\in DD^{-1}$ is a product of at most four original characters or their inverses inflated to $G_M$. It therefore factors through $G_h$ with $h$ the least common multiple of their conductors; $h\mid M$ and $h\le Q^4$. The corresponding character modulo $h$ is nonprincipal because its pullback is nonprincipal. By Lemma~\ref{lem:sample}, the empirical distribution of $S_M$ on $G_M$ obeys \eqref{eq:localbias} with $\eta$ and $|S_M|\ge(1-\eta/8)|S|$. All prime values in $S_M$ agree with those of the inflated character sum.

Theorem~\ref{thm:finite} gives the lower bound on $S_M$, which implies \eqref{eq:mixedlower}. If the full fraction is at most the stated upper threshold, the fraction on $S_M$ is at most $(1+\epsilon)/s$, yielding rigidity on $G_M$. For recovery apply \eqref{eq:mixedlower} to the difference of two sparse mixtures and use the common-good-set argument of Theorem~\ref{thm:code}. Erasures lose $u$ further positions.
\end{proof}

\begin{remark}[Sharpness of the limiting mixed threshold]
Fix $k$, set $s=2k$, and choose a prime $q_0\equiv1\pmod s$ (Dirichlet's theorem suffices). There is a character subgroup of order $s$ modulo $q_0$. Every nonprincipal character modulo the prime $q_0$ is primitive, while the principal term is represented by the conductor-$1$ character. For $Q\ge q_0$ large enough that $X>q_0$, every sampled prime is unramified for this dictionary subfamily. The coset-indicator difference from \eqref{eq:indicator} can be chosen to have weight at most $\lfloor |S|/s\rfloor$, since the $s$ cosets partition $S$. By decreasing the bias for fixed $k$ one approaches the upper obstruction $1/(4k)$ for the fraction of correctable errors. More generally the sharpness statement here concerns a family of moduli where the subgroup exists, not every modulus.
\end{remark}

\section{A less expensive sample for recovery alone}
\begin{theorem}[Near-optimal recovery at quadratic sparsity cost, using \GRH]\label{thm:cheap}
Fix $0<\delta<1/4$. For all sufficiently large $Q$ and all $2\le s\le Q$, there is a universal set of $O_\delta(s^2\log Q)$ distinct primes of size $O_\delta(s^2\log^2 Q)$. On the unramified sample for every modulus $q\le Q$, and on the full sample for the mixed dictionary, every nonzero $s$-sparse sum has $|F(p)|>\sqrt\delta\,\norm c_2$ on at least a fraction
\begin{equation}\label{eq:cheap}
 (1-\delta/(8s))\frac{1-2\delta}{s}
\end{equation}
of the relevant sample. For the fixed-modulus case the prefactor $1-\delta/(8s)$ can be omitted. Hence with $s=2k$, the worst-case correctable error fraction is at least half of \eqref{eq:cheap} (strictly below the half-threshold), and the coefficient error for two compatible noisy candidates is at most $2v/\sqrt\delta$. No extremizer classification is claimed at this coarser bias.
\end{theorem}
\begin{proof}
Take $b=\delta/s$ in Lemma~\ref{lem:sample}; its lower parameter bound holds for large $Q$ depending on $\delta$ and $sb=\delta<1/4$. All fixed-modulus and local mixed-conductor bias conditions hold with $b$ instead of $\eta$. Use only the proof of \eqref{eq:basic}, with $\tau=\sqrt\delta$ and $(t-1)b\le\delta$. The local fraction is at least $(1-2\delta)/s$; in the mixed case the unramified fraction is at least $1-b/8$. The two-candidate proof gives the claimed recovery bound. The sample size and prime cutoff follow from $b^{-2}=s^2/\delta^2$.
\end{proof}

This yields arbitrarily good relative approximation to the optimal uniform error radius with a quadratic, rather than sextic, sparsity cost. The refined sample pays for the additional rigidity and a smaller detection threshold. Neither existence proof requires randomness in the decoder. Randomness is only used to prove that a common distinct prime set exists; a finite search for a successful set with certified character biases is another, potentially very expensive, construction.

\section{Verification, relationship to prior work, and limits}
\label{sec:audit}
\paragraph{Repository analysis.} The merged article's README records 131 reports and 914 results, with the Dirichlet GRH theorem accepted as input, and explicitly says the merged claims were not independently checked \cite{article}. The source section describes the determinant comparison, spanning, torsion selection, and one-sided prime-error bound. In the Lean source, the declaration \texttt{New.grh} supplies nonvanishing to the right of the critical line. It is located in \texttt{OneSidedBound.lean}. The declarations \texttt{New.zeros\_on\_line} and \texttt{New.zeta\_zeros\_on\_line}, both in \texttt{New/Main.lean}, use symmetry to obtain the zero-location statements. We use those outputs only. In particular, we do not attempt to extend the geometric or automorphic part of the proof to other $L$-functions.

Our analytic estimates parallel three labelled results in the merged article: \texttt{dir:lem:smooth} treats principal and real characters; \texttt{rec:lem:gram} is a uniform smoothed pairwise character estimate; and \texttt{dir:lem:oneset} samples distinct primes with a uniform small-bias guarantee at fixed bias. The product and partition results \texttt{dir:thm:quotient}, \texttt{dir:thm:factorization} and \texttt{dir:thm:minority} exploit large character biases to classify product distributions after deletion. Here the direction is different: near-minimum support of a \emph{linear} combination forces the full frequency ratio set into an exact subgroup, even with approximate zeros, and then yields an exact nonlinear code-distance formula. The article's reconstruction results (e.g. \texttt{rec:thm:lowzero}, \texttt{rec:thm:robust}) concern information extracted from zeros or other $L$-function statistics, not the sampled character-mixture code defined here.

\paragraph{External context and novelty.} The Donoho--Stark inequality and exact equality cases on finite abelian groups are classical \cite{MOP}; Meshulam refines the uncertainty inequality in terms of divisors of the group order \cite{Meshulam}. Uncertainty principles in sparse recovery have an extensive literature (e.g. \cite{IM}). The standard Gram estimate, Bernstein sampling, and the Hamming-ball argument are not claimed as new. The proposed contribution is the combination of the local-bias quantitative rigidity theorem (including exact frequency structure under a nonzero tolerance), the exact empirical minimum-distance formula for small primes, and the mixed-conductor transfer with only four-conductor tests. A targeted scan of the relevant article sections and related literature found no matching statement; this is not an exhaustive literature search or a priority certificate. The elementary finite-group argument is independent of GRH and may be useful for other small-bias constructions.

\paragraph{Artifacts and limitations.} The companion \texttt{verify.py} uses exact integer/rational arithmetic to check the constants in the proof, enumerate support distances in the Walsh group of order eight, and produce a balanced-prime example for the quadratic character quotient of modulus $105$. These tests illustrate rather than prove the general theorem. The file \texttt{InputAudit.lean} records the imported theorem interface and is intended for the upstream Lake project; this environment has no Lean installation, so it is not a checked formalization. The mathematical proof above is self-contained apart from classical analytic-number-theory estimates, finite abelian duality and Bernstein's inequality, with the zero-location theorem used as identified. The constants $B,K$ are effective in principle given the classical bounds, but not numerically optimized; we do not claim an efficient decoding algorithm, a deterministic polynomial-time sample construction, or a sharp nondivisibility gap.

\begin{thebibliography}{99}
\bibitem{vendor} ProveIt contributors, \emph{Vendored grh-lean (readable edition)}, repository path \texttt{Analysis/GeneralizedRiemannHypothesis/vendor/grh-lean}, main snapshot accessed 27 Nov. 2026. \url{https://github.com/VladimirReshetnikov/ProveIt/tree/main/Analysis/GeneralizedRiemannHypothesis/vendor/grh-lean}
\bibitem{upstream} L. Alp\"oge, with Claude; L. Sernau and collaborators (rewrite), \emph{The generalized Riemann hypothesis for Dirichlet L-functions, in Lean 4}, commit 72fead102f663d5c7a407964a13dc0cf9d977207, 25 Nov. 2026. \url{https://github.com/lsernau/grh-lean/tree/72fead102f663d5c7a407964a13dc0cf9d977207/grh}
\bibitem{sutherland} A. V. Sutherland, \emph{Verification files for the shorter Lean proof of RH and of GRH for Dirichlet L-functions}, updated 26 Nov. 2026. \url{https://klein.mit.edu/~drew/shortgrh/}
\bibitem{article} ProveIt contributors, \emph{Beyond the Riemann hypothesis for Dirichlet L-functions: consequences, transfers and limits}, 26 Nov. 2026, LaTeX source and README. \url{https://github.com/VladimirReshetnikov/ProveIt/tree/main/Analysis/GeneralizedRiemannHypothesis/Article}
\bibitem{MV} H. L. Montgomery and R. C. Vaughan, \emph{Multiplicative Number Theory I: Classical Theory}, Cambridge University Press, 2007.
\bibitem{Davenport} H. Davenport, \emph{Multiplicative Number Theory}, 3rd ed., revised by H. L. Montgomery, Springer, 2000.
\bibitem{MOP} E. Matusiak, M. \"Ozayd\i n, T. Przebinda, The Donoho--Stark uncertainty principle for a finite abelian group, \emph{Acta Math. Univ. Comenian.} 73 (2004), 155--160.
\bibitem{Meshulam} R. Meshulam, An uncertainty inequality for finite abelian groups, \emph{European J. Combin.} 27 (2006), 63--67. \url{https://arxiv.org/abs/math/0312407}
\bibitem{IM} A. Iosevich and A. Mayeli, Uncertainty principles on finite abelian groups, restriction theory, and applications to sparse signal recovery, preprint (2023). \url{https://arxiv.org/abs/2311.04331}
\end{thebibliography}
\end{document}
